\section{光源空间图集}
\label{sec:method}

\begin{figure*}[t]
  \centering
  \resizebox{0.94\textwidth}{!}{% Method overview. First draft generated with the local Codex CLI (paper/codex/pipeline), then
% revised by hand: layout grid, overlaps, labels. Thumbnails are real buffers of the final Hotdog model.
\begin{tikzpicture}[
  x=1cm,y=1cm,
  font=\sffamily\fontsize{7.4}{8.6}\selectfont,
  line width=0.65pt,
  >={Stealth[length=4pt]},
  light/.style={draw=orange!70!black,fill=orange!8,rounded corners=2pt},
  camera/.style={draw=blue!60!black,fill=blue!6,rounded corners=2pt},
  learned/.style={draw=green!45!black,fill=green!8,rounded corners=2pt},
  lightflow/.style={->,draw=orange!70!black},
  cameraflow/.style={->,draw=blue!60!black},
  mergeflow/.style={->,draw=black!75},
  label/.style={fill=white,inner sep=1.2pt,align=center},
  thumb/.style={draw=gray!65,inner sep=0pt,outer sep=0pt},
  title/.style={font=\sffamily\bfseries\fontsize{7.8}{9}\selectfont},
  stage/.style={align=left,anchor=north west,inner sep=0pt,font=\sffamily\fontsize{7}{8}\selectfont}
]
\path[use as bounding box] (0,0) rectangle (17.4625,6.35);

% ---------------------------------------------------------------- light pass
\node[title,anchor=west,text=orange!70!black] at (0.05,6.18) {光源渲染通道};
\node[light,minimum width=2.40cm,minimum height=1.06cm,align=center] (gaussians) at (1.30,4.62)
  {\textbf{三维高斯}\\[3pt]$\{\mu_i,\Sigma_i,o_i,\mathbf b_i,\mathbf f_i\}$};
\node[learned,minimum width=1.55cm,minimum height=0.43cm] (phi) at (1.30,5.62) {$\boldsymbol\phi$-MLP};
\draw[lightflow] (gaussians.north) -- (phi.south);
\draw[lightflow] (phi.east) -- (2.55,5.62) -- (2.55,6.02) -- (4.90,6.02) -- (4.90,5.88);
\node[label] at (3.72,6.20) {从光源视角光栅化 $\mathbf p$};

\draw[light] (2.94,3.52) rectangle (6.88,5.88);
\node[title] at (4.91,5.70) {光源空间图集};
\node[thumb] at (3.93,4.92) {\includegraphics[width=1.22cm,height=1.22cm]{figures/decomp_hotdog/atlas_depth.png}};
\node[thumb] at (5.86,4.92) {\includegraphics[width=1.22cm,height=1.22cm]{figures/decomp_hotdog/atlas_flux.png}};
\node[align=center] at (3.93,4.08) {深度 $M_1,M_2$};
\node[align=center] at (5.86,4.08) {通量 $\boldsymbol\Phi$};
\node at (4.91,3.72) {$\mathbf A(u)=\sum_i w_i(u)\,[\boldsymbol\phi_i,\,z_i,\,z_i^2,\,1]$};

\draw[lightflow] (6.88,4.92) -- (7.55,4.92);
\node[align=center] at (7.21,4.52) {模糊\\$\downarrow 2$};
\node[title,anchor=south] at (8.50,5.62) {金字塔 $\{\mathbf A_k\}_{k<K}$};
\node[thumb] at (8.25,4.95) {\includegraphics[width=1.22cm,height=1.22cm]{figures/decomp_hotdog/atlas_flux.png}};
\node[thumb] at (8.70,4.66) {\includegraphics[width=0.90cm,height=0.90cm]{figures/decomp_hotdog/atlas_flux_level3.png}};
\node[thumb] at (9.06,4.42) {\includegraphics[width=0.62cm,height=0.62cm]{figures/decomp_hotdog/atlas_flux_level3.png}};

% ---------------------------------------------------------------- camera pass
\node[title,anchor=west,text=blue!60!black] at (0.05,3.38) {相机渲染通道};
\node[camera,minimum width=2.40cm,minimum height=1.20cm,align=center] (receivers) at (1.30,2.42)
  {\textbf{接收点 $\mathbf x$}\\[2pt]高斯中心 (II)\\或像素 (III)};
\node[camera,minimum width=5.90cm,minimum height=1.20cm,inner sep=2pt,align=center] (gather) at (6.02,2.42)
  {\textbf{采样位置 $u(\mathbf x)$, 各层级 $k$}\\[2pt]
   $\mathbf q_k=[M_{0,k},\,t_k,\,\delta_k,\,\log\eta_k,\,V_k],\ \ \boldsymbol\Phi_k$\\[2pt]
   $V_k=1-M_{0,k}\,G(t_k-\beta)$};
\draw[cameraflow] (receivers.north) -- (1.30,3.24) -- (7.90,3.24) -- (7.90,4.32);
\node[label] at (4.85,3.24) {投影至光源视图： $(u(\mathbf x),z(\mathbf x))$};
\draw[cameraflow] (8.70,4.20) -- (8.70,3.03);
\node[label,anchor=west] at (8.78,3.62) {采样};

% ---------------------------------------------------------------- learned heads
\node[title,anchor=west,text=green!45!black] at (10.05,6.18) {着色};
\node[learned,minimum width=4.85cm,minimum height=1.05cm,align=center] (visibility) at (12.50,5.42)
  {\textbf{可见性 $\mathcal G$}\\[2pt]$V=\sigma\big(\mathrm{logit}\,V_0+\mathcal G(\mathbf Q,\mathbf f,\mathbf n\!\cdot\!\mathbf l)\big)$};
\node[learned,minimum width=4.85cm,minimum height=1.05cm,align=center] (transfer) at (12.50,4.14)
  {\textbf{光传输 $\mathcal K$}\\[2pt]$L_{\mathrm{tr}}=\sum_k\mathbf W_k\,\boldsymbol\Phi_k$, \ $\mathbf W_k$ 来自 $\mathbf Q$};
\node[learned,minimum width=4.85cm,minimum height=1.05cm,align=center] (material) at (12.50,2.86)
  {\textbf{局部材质}\\[2pt]$\rho=\mathrm{MLP}(\mathbf f,\mathbf n,\mathbf l,\mathbf v,\mathbf h,\mathbf r)$, 瓣函数 $s$};
\draw[draw=blue!60!black] (gather.east) -- (9.70,2.42) -- (9.70,5.42);
\draw[cameraflow] (9.70,5.42) -- (visibility.west);
\draw[cameraflow] (9.70,4.14) -- (transfer.west);
\draw[cameraflow] (9.70,2.86) -- (material.west);
\fill[blue!60!black] (9.70,4.14) circle (0.65pt);
\fill[blue!60!black] (9.70,2.86) circle (0.65pt);

\node[thumb] (visimage) at (16.20,5.42) {\includegraphics[width=1.12cm,height=1.12cm]{figures/decomp_hotdog/visibility.png}};
\node[thumb] (trimage) at (16.20,4.14) {\includegraphics[width=1.12cm,height=1.12cm]{figures/decomp_hotdog/transfer.png}};
\node[thumb] (matimage) at (16.20,2.86) {\includegraphics[width=1.12cm,height=1.12cm]{figures/decomp_hotdog/local.png}};
\draw[draw=green!45!black] (visibility.east) -- (visimage.west);
\draw[draw=green!45!black] (transfer.east) -- (trimage.west);
\draw[draw=green!45!black] (material.east) -- (matimage.west);

% ---------------------------------------------------------------- radiance
\draw[draw=black!75] (15.25,5.42) -- (15.25,2.03);
\fill[black!75] (15.25,5.42) circle (0.65pt);
\fill[black!75] (15.25,4.14) circle (0.65pt);
\fill[black!75] (15.25,2.86) circle (0.65pt);
\node[draw=black,fill=white,rounded corners=2pt,minimum width=15.0cm,minimum height=0.58cm,align=center] (shading) at (7.62,1.45)
  {$L(\mathbf x)=\bar E(\mathbf x)\big[\,V\rho+V_0\,s+\textstyle\sum_k\mathbf W_k\boldsymbol\Phi_k\big],\qquad \bar E(\mathbf x)=\mathbf I/\|\mathbf p-\mathbf x\|^2$};
\draw[mergeflow] (15.25,2.03) -- (15.25,1.74);
\node[thumb] (output) at (16.55,1.45) {\includegraphics[width=1.05cm,height=1.05cm]{figures/decomp_hotdog/ours.png}};
\draw[mergeflow] (shading.east) -- (output.west);

% ---------------------------------------------------------------- schedule (8k : 22k : 8k)
\fill[gray!10] (0.06,0.02) rectangle (3.7116,0.98);
\fill[blue!8] (3.7116,0.02) rectangle (13.7534,0.98);
\fill[green!8] (13.7534,0.02) rectangle (17.405,0.98);
\draw[draw=gray!60,rounded corners=2pt] (0.06,0.02) rectangle (17.405,0.98);
\draw[draw=gray!60] (3.7116,0.02) -- (3.7116,0.98);
\draw[draw=gray!60] (13.7534,0.02) -- (13.7534,0.98);
\node[stage,text width=3.50cm] at (0.16,0.90)
  {\textbf{I. 几何预热} (8k)\\无网络；目标 $Y^{1/\gamma}\!\to\!Y$};
\node[stage,text width=9.80cm] at (3.84,0.90)
  {\textbf{II. 联合训练} (22k)\\线性辐亮度损失；高斯接收点；增密};
\node[stage,text width=3.55cm] at (13.86,0.90)
  {\textbf{III. 精修} (8k)\\冻结几何；像素接收点；显示域损失};
\end{tikzpicture}
}
  \caption{\textbf{方法概览。}对每个光源，从光源视角光栅化高斯一次；每个高斯按其截获光比例对应的合成权重，沉积可学习通量特征 $\boldsymbol\phi_i$ 与深度矩（\cref{eq:atlas}）。接收点投影到光源视图，从金字塔汇集多尺度统计量与通量，并结合矩阴影检验上的残差、线性作用于汇集通量特征的光传输核，以及局部材质进行着色。缩略图：Hotdog 模型在测试视图上的实际缓冲区。底部：\cref{sec:optimization} 的训练流程。}
  \label{fig:pipeline}
\end{figure*}


我们从图像集合 $\{Y_m\}$ 重建物体。每张图像均由经过标定的相机拍摄，并由位于 $\vect p_m$、RGB 强度为 $\vect I_m$ 的点光源照明（\cref{fig:pipeline}）。

\subsection{预备知识}
\label{sec:prelim}
场景由一组三维高斯组成，其均值、协方差和不透明度分别为 $\mu_i$、$\Sigma_i$ 和 $o_i$~\cite{kerbl20233d}。
每个高斯携带基础颜色 $\vect b_i\in\mathbb R^3$ 和潜在材质编码 $\vect f_i\in\mathbb R^{32}$，其中前三个通道定义可学习的单位法线 $\vect n_i$。
泼溅按从前到后的顺序合成逐高斯属性 $\vect c_i$：
\begin{equation}
  \bar{\vect c}(u)=\sum_i w_i(u)\,\vect c_i,\qquad w_i(u)=\alpha_i(u)\prod_{j<i}\big(1-\alpha_j(u)\big),
  \label{eq:composite}
\end{equation}
其中，$\alpha_i(u)$ 为高斯 $i$ 在像素 $u$ 处的不透明度。
对于点 $\vect x$，以 $\vect l,\vect v,\vect h$ 分别表示单位光照方向、视线方向和半程向量，并以 $\bar E(\vect x)=\vect I/\|\vect p-\vect x\|^2$ 表示法向入射时的辐照度。

\subsection{从光源视角泼溅以记录沉积量}
\label{sec:deposit}
在光源位置放置一台朝向场景中心的透视相机，视场覆盖 99.9\% 高斯的泼溅范围。
其像素 $u$ 对应立体角 $d\omega_u$，接收通量 $\vect I\,d\omega_u$。
在泼溅模型中，每个高斯吸收沿射线继续传播的光中 $\alpha_i(u)$ 的比例，因此其截获通量为
\begin{equation}
  P_i(u)=\vect I\,d\omega_u\,w_i(u),
  \label{eq:flux}
\end{equation}
其中权重由 \cref{eq:composite} 给出。
这一关系成立于光栅化泼溅模型内部；该模型通过瓦片内按中心深度排序、低通滤波和不透明度截断，近似体积透射率~\cite{radl2024stopthepop,huang2024error,talegaonkar2025volumetrically,yu2024mip}。
\cref{eq:flux} 不显式包含法线或距离：距离为 $r$、入射角为 $\theta$ 的面元 $dA$ 接收 $\vect I\cos\theta\,dA/r^2=\vect I\,d\omega$，因此投影缩短和距离衰减已经由表面覆盖的光源像素数量体现。

\paragraph{图集。}
每个高斯输出 $\vect c_i=[\boldsymbol\phi_i,\,z_i,\,z_i^2,\,1]$，其中 $z_i$ 是其中心沿光源光轴的深度，以场景中心为基准并用场景半径归一化；$\boldsymbol\phi_i=\softplus(\mathrm{MLP}_\phi(\vect f_i,\vect b_i,\vect n_i\cdot\vect l_i))\in\mathbb R_+^{C}$ 为 $C=8$ 个可学习通量特征。
一次 $512^2$ 光栅化得到图集
\begin{equation}
  \vect A(u)=\sum_i w_i(u)\,\vect c_i=\big[\boldsymbol\Phi(u),\,M_1(u),\,M_2(u),\,M_0(u)\big],
  \label{eq:atlas}
\end{equation}
它表示以该像素入射通量为单位的沉积量：$M_0$ 是被场景吸收的通量比例，$M_1/M_0$ 和 $M_2/M_0$ 是吸收深度的前两个矩，$\boldsymbol\Phi$ 则描述光落在了何种物体上。
绝对沉积量 $\vect I\,d\omega_u\vect A(u)$ 还涉及立体角因子，且 $d\omega_u\propto\cos^3$，其中余弦对应偏离光轴的角度。
在我们的数据中，典型图集内该量的变化小于 30\%（最宽视场的图集角落可相差 $2.1\times$）；下文的阴影统计量是比值，不依赖它，而光传输项忽略了这一因子。
反复进行二项式模糊和 $2\times$ 下采样，得到具有 $K=7$ 个层级的金字塔 $\vect A_0,\dots,\vect A_{K-1}$。
由于每个通道对权重 $w_i$ 都呈线性，各层仍保存相应覆盖区域内沉积量的矩与特征和；覆盖范围从单个纹素逐渐增至约三分之一物体大小，这正是此类阴影图可以进行滤波的原因~\cite{donnelly2006variance,peters2015moment}。

\subsection{根据图集着色}
\label{sec:shading}
接收点 $\vect x$ 以归一化深度 $z(\vect x)$ 投影至光源视图的 $u(\vect x)$，并对每个层级进行双线性采样。
在层级 $k$，我们构造平均沉积深度 $\bar z_k=M_{1,k}/M_{0,k}$、包含纹素尺寸 $\tau_k$ 的离散程度 $\eta_k=(M_{2,k}/M_{0,k}-\bar z_k^2+\tau_k^2)^{1/2}$、接收点相对沉积深度的后向偏移 $\delta_k=z-\bar z_k$，以及 $t_k=\delta_k/\eta_k$。
若覆盖区域内的沉积深度服从正态分布，则在接收点之前吸收的光比例为 $M_{0,k}\,G(t_k)$，其中 $G$ 为标准正态累积分布函数，由此得到基于矩的检验：
\begin{equation}
  V_k=1-M_{0,k}\,G(t_k-\beta),
  \label{eq:moment}
\end{equation}
其中偏置 $\beta=3$ 以局部离散程度为单位，用于避免自阴影。
所有层级的描述子 $\vect q_k=[M_{0,k},t_k,\delta_k,\log\eta_k,V_k]$ 组成 $\vect Q(\vect x)\in\mathbb R^{5K}$。
当遮挡物与接收点处于同一覆盖区域时，这一检验只能提供较弱的先验（例如，不透明接收点前方有半透明遮挡物时，$t_0=1$，\cref{eq:moment} 返回的是 $0.98$ 而非 $0.5$，方差阴影图的切比雪夫界则能给出正确结果~\cite{donnelly2006variance}）。
此外，每个高斯存入的是中心深度，因此一个受到斜向照明的大型高斯会被表示为单一深度。

\paragraph{可见性。}
以零初始化的残差网络 $\mathcal G$ 修正最细层级的检验：
\begin{equation}
  V(\vect x)=\sigma\big(\logit V_0(\vect x)+\mathcal G(\vect Q,\vect f,\vect n\cdot\vect l)\big),
  \label{eq:visibility}
\end{equation}
因此，训练从经典滤波阴影图出发，$\mathcal G$ 根据多尺度统计量学习修正上述情形，以及毛发的部分透光和真实有限尺寸光源产生的柔软阴影边界。

\paragraph{光传输。}
根据 \cref{eq:flux}，进入半透明物体并在 $\vect x$ 处离开的光，其贡献为 $\vect I\!\int\!\sum_i R(\vect x,\mu_i)\,w_i(\omega)\,d\omega$：即将依赖接收点的核 $R$ 与沉积量相乘，在光源立体角范围内积分。
半透明阴影图通过存储辐照度和固定扩散核计算此类积分~\cite{dachsbacher2003translucent,deon2007efficient,jimenez2010real}。
我们同时学习沉积特征和核，并将金字塔层级作为宽度递增的核：
\begin{align}
  L_{\mathrm{tr}}(\vect x)&=\textstyle\sum_{k}\vect W_k\,\boldsymbol\Phi_k(\vect x),\label{eq:transfer}\\
  \{\vect W_k\}_k&=\softplus\big(\mathcal K(\vect Q,\vect f,\vect n\cdot\vect l,\vect n\cdot\vect v)\big),\quad \vect W_k\in\mathbb R_+^{3\times C}.\nonumber
\end{align}
这些核仅依赖接收点及几何统计量 $\vect Q$，因此在给定几何与光源时，$L_{\mathrm{tr}}$ 对沉积特征呈线性。
由于接收点自身的沉积量也位于其覆盖区域内，$L_{\mathrm{tr}}$ 同样可能吸收部分直接反射，各项之间的分配由学习决定。

\paragraph{局部材质与辐亮度。}
局部反射率 $\rho(\vect x)$ 是由潜在编码驱动的神经 BRDF：MLP 的输入包括 $\vect f$、$\vect n$、$\vect l,\vect v,\vect h$ 及反射视线方向 $\vect r$ 的傅里叶编码、相关余弦项和高光提示 $e^{\kappa(\vect n\cdot\vect h-1)}$~\cite{zeng2023nrhints}，其输出再乘以平滑余弦 $\psi(\vect n\cdot\vect l)=\softplus(20\,\vect n\cdot\vect l)/20$。
位置信息仅通过八个与光源无关的系数进入模型，用于混合八个依赖光源的响应，从而限制投射阴影被固化到材质中的程度（自由的逐高斯编码仍允许其中一部分）。
各向同性瓣函数 $s(\vect x)$ 采用固定锐度 $\kappa\in\{8,\dots,2048\}$ 和与光源无关的权重，负责表达尖锐高光。
出射辐亮度为
\begin{equation}
  L(\vect x)=\bar E(\vect x)\big[V(\vect x)\,\rho(\vect x)+V_0(\vect x)\,s(\vect x)+L_{\mathrm{tr}}(\vect x)\big].
  \label{eq:radiance}
\end{equation}
瓣函数使用几何检验 $V_0$，使高光在可学习的 $V$ 尚不准确时仍保有梯度传播路径。
将 $\bar E(\vect x)$ 提取到括号外，使网络能够预测类似反射率的量；对 $L_{\mathrm{tr}}$ 而言，这相当于以接收点处的距离衰减替代入射点处的衰减，在背光透射情形中会产生两者到光源距离平方比 $(r_{\mathrm{out}}/r_{\mathrm{in}})^2$ 的偏差。
核只能在平均意义上吸收这一误差；当同一次拍摄中的光源距离变化较小时，这种近似足够有效（合成场景中距离不变，真实场景中变化不超过 $1.8\times$；见补充材料）。
$\mathcal G$ 和 $\mathcal K$ 的所有输入都相对于光源定义，因此算子定义在光源坐标系中，而非按光源位置索引；这是否有助于外推，将由实验 E1 检验（\cref{sec:exp_atlas}）。

\subsection{接收点：高斯或像素}
\label{sec:receivers}
\emph{高斯接收点}对每个高斯中心 $\mu_i$ 着色，再用 \cref{eq:composite} 合成颜色；\emph{像素接收点}则先合成 $\vect b,\vect f$ 和期望深度，将每个像素反投影为世界坐标点，再进行一次着色。
高斯接收点为几何提供增密所依赖的常规光栅化梯度，并避免重叠表面的材质混合，但也继承了 \cref{sec:intro} 所述逐图元阴影的局限。
像素接收点在每个像素对应的表面位置求值阴影检验，因此阴影边界可以穿过大型高斯。
与 RNG 从前向着色切换到延迟着色的做法类似~\cite{fan2025rng}，\ours 在几何优化期间使用高斯接收点，几何固定后改用像素接收点（\cref{sec:optimization}），两者共享同一图集。
